\section{Experiment 9: Mechanism Audit}
\label{sec:exp9}

\subsection{Purpose}

Experiments~1--8 establish that finite-memory event communication can strongly reduce communication while retaining good scalar optimization accuracy. They do \emph{not}, however, establish that the LIF mechanism is algorithmically distinct from standard filtering, error-feedback, or event-thresholding constructions. Experiment~9 is therefore a gatekeeping experiment. Its question is
\begin{equation}
  \boxed{\text{Does the current FedLIF mechanism provide a distinct error--communication trade-off?}}
\end{equation}
The goal is explicitly not to obtain another favorable operating point. Each comparator is tuned over a small parameter grid and evaluated through an error--communication Pareto front.

\subsection{Reference problem and fairness conventions}

The audit keeps the heterogeneous delayed-asynchronous scalar problem from Experiment~8. Two equally weighted clients optimize
\begin{equation}
  F_i(w)=\frac12(w-\theta_i)^2,
  \qquad
  \theta_{\mathrm{slow}}=+0.05,
  \quad
  \theta_{\mathrm{fast}}=-0.05,
\end{equation}
so the intended aggregate optimum is $w^\star=0$. The slow/fast compute-delay ratio is $R=40$. A client computes on a server snapshot and returns the corresponding stochastic gradient only after its compute period. Compute-rate normalization is retained:
\begin{equation}
  \gamma_i=\gamma p_iT_i,
\end{equation}
so faster hardware does not receive larger optimization weight solely from producing more gradient results per wall-clock time.

The primary noisy setting uses $\sigma=0.25$. A deterministic $\sigma=0$ control is retained. Initial conditions are randomized over a narrow interval to reduce lattice artifacts. All methods at a fixed noise setting use common random numbers. Event-based methods are tuned over the event quantum $q$ and their communication/filter parameters. Full-precision and periodic-sign anchors are also tuned over their step size or event quantum rather than inheriting one arbitrary setting.

In this scalar experiment, communication cost is represented by the number of transmitted signed events/messages. This is only a mechanism-level proxy. In the vector experiments, actual bit accounting must include coordinate or block addresses, sign bits, packetization, and protocol overhead.

\subsection{Comparator families}

The following mechanisms are compared.
\begin{enumerate}
  \item \textbf{FedLIF with subtractive reset.} The current working algorithm,
  \begin{align}
    z_i^+ &= \rho z_i + u_i,\\
    |z_i|\geq\Delta &\Rightarrow
    \begin{cases}
      w^+ = w + q\,\sign(z_i),\\
      z_i^+ = z_i-\Delta\sign(z_i).
    \end{cases}
  \end{align}
  \item \textbf{Full-reset LIF.} The same leaky integration and threshold, but a firing event resets the membrane to rest, $z_i^+=0$.
  \item \textbf{IF with remote reset.} Infinite local memory, but a remote-client event invalidates the other client's accumulator. This is a simple explicit freshness mechanism.
  \item \textbf{Memoryless pulse quantization.} The current rate-normalized gradient evidence is quantized directly into signed pulses without temporal memory.
  \item \textbf{EMA plus memoryless pulse quantization.} A conventional exponential gradient filter is applied before stateless pulse quantization. This separates low-pass filtering from persistent residual accumulation.
  \item \textbf{Periodic sign.} A fixed signed update is sent for every completed stochastic gradient.
  \item \textbf{Full-precision asynchronous gradient updates.} This provides an accuracy anchor at high communication cost.
\end{enumerate}

\subsection{Exact equivalence results}

The audit reveals an important novelty boundary that is stronger than any empirical comparison.

\paragraph{Subtractive FedLIF is decayed error feedback after rescaling.}
Let $u_k$ denote the LIF evidence increment and write the subtractive-reset recurrence abstractly as
\begin{equation}
  z_{k+1}=\rho z_k+u_k-\Delta s_k,
\end{equation}
where $s_k\in\{-1,0,+1\}$ denotes emitted threshold events. Define a residual in parameter-update units
\begin{equation}
  e_k=\frac{q}{\Delta}z_k.
\end{equation}
Then
\begin{equation}
  e_{k+1}=\rho e_k+\frac{q}{\Delta}u_k-q s_k.
\end{equation}
This is precisely a decayed error-feedback/residual recurrence with a fixed signed quantizer. The numerical equivalence check gives zero mismatching events/model states up to floating-point roundoff. Error feedback with memory decay already appears in communication-efficient learning, e.g. CO$^3$~\cite{chen2022co3}, while conventional gradient sparsification also accumulates unsent information~\cite{lin2017dgc,stich2018memory}. Consequently, the subtractive FedLIF recurrence alone cannot be claimed as a novel optimizer mechanism.

\paragraph{Full-reset LIF is a reset EMA trigger after rescaling.}
For one client with evidence gain $c_i$, full-reset LIF obeys
\begin{equation}
  z_{k+1}=\rho z_k-c_i g_k.
\end{equation}
Define
\begin{equation}
  m_k=-\frac{1-\rho}{c_i}z_k.
\end{equation}
Between reset events,
\begin{equation}
  m_{k+1}=\rho m_k+(1-\rho)g_k,
\end{equation}
which is the standard exponential moving average. The firing threshold maps to
\begin{equation}
  |m_k|\geq \frac{(1-\rho)\Delta}{c_i},
\end{equation}
and the LIF reset $z^+=0$ is exactly $m^+=0$. The numerical equivalence check again yields zero event/model mismatches up to floating-point roundoff. Thus the full-reset form is a neuromorphic interpretation of a thresholded resetting EMA, not a fundamentally different recurrence.

These identities are central conclusions of Experiment~9. The project may still obtain value from the hybrid/event-driven interpretation, temporal coding, or additional neuromorphic design rules, but novelty cannot rest on either base recurrence alone.

\subsection{Tuned noisy result}

A higher-seed confirmation run at $R=40$ and $\sigma=0.25$ gives the representative best configurations in \cref{tab:exp9-noisy}.

\begin{table}[ht]
\centering
\caption{Experiment~9 tuned comparator summary in the noisy heterogeneous delayed-asynchronous setting. Event/message counts are scalar communication proxies, not complete bit costs.}
\label{tab:exp9-noisy}
\begin{adjustbox}{max width=\textwidth}
\begin{tabular}{lrrr}
\toprule
Method & Mean events/messages & Tail RMSE & Harmful-event fraction\\
\midrule
Full precision & 4612.0 & 0.02833 & 0.375\\
Full-reset LIF & 28.33 & 0.03174 & 0.146\\
FedLIF, subtractive reset & 28.88 & 0.03386 & 0.154\\
IF & 31.38 & 0.03602 & 0.183\\
IF with remote reset & 27.71 & 0.04129 & 0.142\\
EMA + memoryless pulse & 42.89 & 0.08873 & 0.278\\
Periodic sign & 4612.0 & 0.10428 & 0.498\\
Memoryless pulse & 164.75 & 0.12885 & 0.445\\
\bottomrule
\end{tabular}
\end{adjustbox}
\end{table}

The event-driven finite-memory methods therefore remain highly effective as communication mechanisms in this toy problem. Full-reset LIF reaches a tail RMSE only slightly above the tuned full-precision reference while replacing thousands of messages by a few tens of signed events. It also marginally outperforms the subtractive-residual variant in this setting. This suggests that preserving every threshold overshoot is not essential here and can even retain undesirable noisy/heterogeneous evidence.

However, the empirical advantage over the deliberately simpler EMA-plus-memoryless-pulse baseline does \emph{not} establish novelty over EMA filtering in general, because the exact rescaling above shows that a thresholded \emph{resetting} EMA can reproduce full-reset LIF exactly. The comparator family must therefore be interpreted together with the equivalence result, not in isolation.

\subsection{Deterministic control}

In the $\sigma=0$ control, IF, subtractive FedLIF, and full-reset LIF all obtain tail errors of order $2\times10^{-2}$ with roughly $25$--$32$ events for their best tuned configurations, while full precision reaches a much smaller numerical neighborhood at high communication cost. The relative finite-memory advantage is therefore markedly stronger under stochastic gradients, consistent with the earlier conclusion that temporal noise filtering is one of the robust mechanisms.

\subsection{Interpretation and decision}

Experiment~9 produces a mixed but decisive result.
\begin{enumerate}
  \item \textbf{The event-driven temporal mechanism is useful.} Properly tuned finite memory achieves an excellent scalar error--communication trade-off and is substantially better than periodic sign or memoryless pulse baselines.
  \item \textbf{The current FedLIF-v0 recurrence is not algorithmically unique.} Subtractive reset is exactly decayed error feedback after rescaling.
  \item \textbf{The empirically strongest full-reset variant is also not unique as a recurrence.} It is exactly a thresholded resetting EMA after rescaling.
  \item \textbf{Explicit remote reset is competitive but weaker.} This reinforces the earlier conclusion that generic stale-gradient invalidation is not the main benefit.
\end{enumerate}

The original Experiment~9 gate should therefore be regarded as only partially passed. There is enough evidence to continue developing \emph{neuromorphic event-driven federation}, but there is not yet enough evidence to scale the present FedLIF-v0 recurrence directly and claim a new optimizer. Before or alongside the vector extension, the project should introduce and test a mechanism that is genuinely tied to the neuromorphic/event-time viewpoint, such as adaptive/homeostatic firing thresholds, firing-rate regulation, refractory dynamics, event-time-aware aggregation, or another hybrid design rule that is not algebraically reducible to standard EMA/error-feedback compression.
